\begin{table}[!htbp]
\centering
\footnotesize
\setlength{\tabcolsep}{3.5pt}
\caption*{\textbf{Table C30.} Conservative reporting gains with primary suppression}
\begin{tabularx}{\linewidth}{lCCCCC}
\toprule
Election & Current width & Exact counts & Suppress counts <5 & Suppress counts <10 & Suppress counts <20 \\
\midrule
2022 & 0.249 & 26.9\% & 26.2\% & 25.3\% & 23.7\% \\
2024 & 0.205 & 29.7\% & 29.0\% & 28.3\% & 26.8\% \\
2025 & 0.231 & 31.9\% & 31.4\% & 30.7\% & 29.6\% \\
2026 & 0.101 & 57.1\% & 55.3\% & 53.2\% & 49.0\% \\
\bottomrule
\end{tabularx}
\par\vspace{2pt}{\raggedright\footnotesize\textit{Notes.} Main assumptions, within-district z slope. Out-of-zone early counts by registered neighborhood and candidate, including invalid ballots; other pooled modes stay pooled. Exact counts reproduce Table 5. Other columns express guaranteed reductions relative to the current continuous range, using an upper bound on remaining integer uncertainty under a release that shows counts >=k exactly and marks smaller counts suppressed. A continuous paired-allocation relaxation allows each count to differ by up to k-1 and can overstate the maximum remaining width. These are not sharp suppression bounds or fractions of an exactly solved integer prior range. Complementary suppression of larger cells and a full privacy assessment are not covered (Appendix A.3).\par}
\end{table}
